\chapter{Problem Description}

\section{The Linear Ordering Problem}

The Linear Ordering Problem (LOP) is a combinatorial optimisation problem
that asks for a permutation of the rows and columns of a square matrix that
maximises the sum of superdiagonal entries. It is $\mathcal{NP}$-hard in
general, which means no polynomial-time exact algorithm is known for it.

More formally, given an $n \times n$ matrix $C = (c_{ij})$ with non-negative
integer entries, the goal is to find a permutation $\pi$ of $\{1, \ldots, n\}$
that maximises
\begin{equation}
    f(\pi) = \sum_{i=1}^{n} \sum_{j=i+1}^{n} c_{\pi_i \pi_j},
    \label{eq:lop_obj}
\end{equation}
where $\pi_i$ is the element placed at position $i$~\cite{SchStu04}. Summing
$c_{\pi_i \pi_j}$ for all $i < j$ collects all matrix entries that end up
above the diagonal once rows and columns are reordered according to $\pi$.

\section{Equivalent Formulations}

The LOP is also known as the \textit{triangulation problem}: permute the rows
and columns of $C$ simultaneously so that the upper triangle is as large as
possible (or equivalently, the lower triangle as small as possible). The two
formulations are equivalent.

There is also a graph interpretation. Given a complete weighted directed graph
on $n$ nodes, a tournament selects one arc from each pair of nodes. An acyclic
tournament defines a total ordering of the nodes, and maximising the total arc
weight of the compatible arcs is exactly the LOP~\cite{MarReiDua12}.

Before running any algorithm, each instance is converted to \textit{normal
form}: all entries are made non-negative, diagonal entries are set to zero,
and for each pair $(i, j)$ the smaller of $\{c_{ij}, c_{ji}\}$ is set to
zero. This transformation does not change the optimal permutation but makes
comparing algorithm results across instances more meaningful.

\section{Applications}

Despite being a theoretical combinatorial problem, the LOP appears in several
practical contexts:
\begin{itemize}
    \item \textbf{Economics:} triangularising input-output tables to reveal the
    hierarchical structure of production sectors, as originally proposed by
    Chenery and Watanabe~\cite{CheWat58}.
    \item \textbf{Archaeology:} finding the most likely chronological ordering
    of artefacts found at an excavation site.
    \item \textbf{Scheduling:} ordering tasks or jobs to minimise
    precedence-constraint violations.
    \item \textbf{Graph drawing:} minimising the number of arc crossings in
    layered directed graphs.
\end{itemize}

\section{Benchmark Instances}

The experiments in this report use a set of 78 benchmark instances provided
for this exercise, with matrix sizes $n = 150$ and $n = 250$ (49 instances of
each size). These instances are large enough to make exact optimisation
intractable, which is precisely what motivates using heuristic approaches.
For each instance, a best-known objective value is available and is used
throughout to compute the relative percentage deviation (RPD), allowing a
fair comparison of solution quality across algorithms.
